\section{Related work}
\label{sec:related}

\subsection{Deep learning for RUL prediction on C-MAPSS}

C-MAPSS \cite{saxena2008damage} simulates turbofan degradation under one or
six operating conditions and one or two fault modes. Early deep models used a
single family: DCNN \cite{babu2016deep}, multi-scale CNNs
\cite{li2018remaining}, LSTMs and GRUs
\cite{zheng2017long,yuan2016fault,xiang2021multicellular}, often with
attention re-weighting \cite{zhang2022prediction,cnnlstmsam2019li}. A second
generation adopted Transformers \cite{vaswani2017attention} through
dimension-independent embeddings \cite{wang2024dvgtformer}, temporal-feature
attention \cite{zhang2024tatfa}, probabilistic heads
\cite{vaetransformer2022} or domain-knowledge augmentation
\cite{kim2025dkam}. Recent applied work mostly uses the hybrid backbone
studied here (TCN + BiGRU + Transformer encoder). RMSE, and occasionally the
PHM08 score \cite{saxena2008damage}, comes from a single run, with seed
variation treated as noise rather than as something worth measuring. Tuning
each subset separately is standard, which can hide a lucky seed on the
multi-condition subsets. We have not found a published failure rate for any
architecture.

\subsection{Handling operating conditions}

On FD002/FD004 the same state of wear produces different sensor readings
across six flight regimes. Proposed remedies include clustering by condition
and normalising each cluster separately \cite{liu2021clustering}, learned
condition embeddings, and training objectives meant to remove condition
information from the internal features, such as gradient reversal
\cite{ganin2016domain,ragab2022contrastive} and contrastive alignment across
conditions. Reported results are mixed and dataset-specific. Our archive
suggests a reason: gradient reversal moves the failing settings around
rather than removing them, by pushing condition information into the very
features most likely to be picked up, while contrastive losses break down
when same-state pairs from different conditions are rare
(Sections~\ref{sec:phenomenon} and~\ref{sec:method_prune}). Recent
multi-condition architectures \cite{kim2025dkam} gain robustness from
multi-scale machinery instead of trying to remove condition information,
which agrees with what we find: keeping more than one branch that can carry
the signal matters more than erasing condition cues.

\subsection{Training failures and attention degeneration}

Work on two other kinds of training failure has not, as far as we know, been
applied to RUL prediction. Shortcut learning
\cite{geirhos2020shortcut,sturm2012clever} is the name for a classifier using
an easy cue in the data instead of the one the task needs; our probes show
the same thing with operating condition (Section~\ref{sec:mechanism}), with
one difference. In regression that cue cannot supply its own target: which
condition an engine is in says nothing about its RUL beyond the average, so
the network settles on the average. Attention degeneration is documented in
vision Transformers, where attention sharpens too early and learning freezes
\cite{zhai2023stabilizing,dong2021attention}. Our case goes the other way:
attention flattens until every time step is weighted equally (entropy 1.0),
so the attention block becomes a fixed average and positional information
disappears (Section~\ref{sec:mechanism}). Feature collapse in self-supervised
learning \cite{chen2021exploring,grill2020bootstrap} and mode collapse in GANs
are related degenerate solutions. None of this work has looked at a hybrid
backbone on multivariate sensor regression, where the outcome also depends
on the data structure and on parallel branches that can hide or rescue the
failure.
